\documentclass[11pt]{scrartcl}
\usepackage[utf8]{inputenc}
\usepackage[sexy]{evanchen} 
\usepackage{graphicx}
\graphicspath{{8 Outsource K Math/Gambar/}}

\title{Digit}
\author{Azzam and Akeyla}
\date{September 2022}

\begin{document}

\maketitle

\section{Properti}
Definisikan $S(n)$ sebagai jumlah digit dari $n$.
\begin{enumerate}
    \item Untuk sembarang bilangan asli $a$ dan $n$ berlaku $S(an) \leq aS(n)$
    \item $S(n) \equiv n \pmod{9}$
    \item $S(ab) \leq S(a)S(b)$
\end{enumerate}


\section{Soal-soal}
    \begin{problem}
    Buktikan properti yang sudah diberikan
    \end{problem}

    \begin{problem}
    Apakah terdapat bilangan asli $N$ sedemikian hingga $N > 2024$ dan ada sebuah bilangan asli $k$ dengan $N$ digit, $S(k) = N$, dan $k$ habis dibagi oleh $N$?
    \end{problem}

    \begin{problem}[IGMO 2022]
    Ambil sebuah bilangan asli $n_0$. Dapatkan $n_1$ sebagai jumlah kuadrat digit-digit $n_0$. Dengan cara yang serupa, konstruksikan $n_2, n_3, $ dan seterusnya. Buktikan bahwa barisan $\{n_i\}_{i = 0}^{\infty}$ selalu suatu saat akan periodik, berapapun nilai $n_0$.
    \end{problem}
    
    \begin{problem}[Rusia MO 1999]
    Dalam representasi desimal dari $N$, digitnya tampil secara menaik. Tentukan semua nilai yang mungkin dari $S(9N)$
    \end{problem}
    
    \begin{problem}[Romania TST JBMO 1999]
    Tentukan semua bilangan asli $n$ sedemikian hingga terdapat bilangan cacah $a$ dan $b$ sedemikian hingga \[S(a) = S(b) = S(a + b) = n\]
    \end{problem}
    
    \begin{problem}[Rusia MO 1997]
    Buktikan bahwa terdapat tak hingga banyaknya bilangan asli $n$ sedemikian hingga $S(3^n) \geq S(3^{n + 1})$
    \end{problem}
    
    \begin{problem}[Rusia MO 1999]
    Diketahui bahwa $S(N) = 100$ dan $S(44N) = 800$. Tentukan nilai dari $S(3N)$
    \end{problem}
    
    \begin{problem}[IMO 1968]
    Tentukan semua bilangan asli $n$ yang hasil kali digit-digitnya sama dengan $n^2 - 10n - 22$
    \end{problem}
    
    \begin{problem}[Czech-Polish-Slovak Junior Match 2017]
    Tentukan bilangan asli $n \geq 3$ terbesar sedemikian hingga terdapat bilangan asli $n$ digit $\bar{a_1a_2...a_n}$ yang nilai $a_1, a_2,$ dan $a_n$ taknol yang habis dibagi oleh $\bar{a_2...a_n}$
    \end{problem}

    \begin{problem}[KTOM April 2022]
        Sebuah bilangan asli $n$ disebut sebagai \textit{Komi} apabila tidak ada digitnya yang bernilai $0$, jumlah digit $n$ sama dengan $2021$, tetapi jumlah digit dari $2n$ tidak lebih dari $1202$. Jika bilangan \textit{Komi} terbesar adalah $N$, banyak digit dari $N$ adalah ...
    \end{problem}
    
\end{document}